\documentclass[11pt]{article}
\usepackage[margin=1in]{geometry}
\usepackage{amsmath,amssymb,amsthm}
\usepackage[hidelinks]{hyperref}
\newtheorem{proposition}{Proposition}
\newcommand{\cSM}{c^{\mathrm{SM}}}
\title{Failure of stable CSM positivity at the empty partition\\
\large MathDB 374937}
\author{}
\date{September 12, 2026}
\begin{document}
\maketitle
\vspace{-2em}

\paragraph{Task 1.}
Resolve the shifted-binomial positivity and log-concavity assertion for
stable CSM--Schur coefficients in \cite[Conjecture 6.7]{BF}.
For a partition $\lambda$ whose parts are at least $2$, put
$c=|\lambda|-\ell(\lambda)$ and write
\[
 \cSM(Y_\lambda(d))_f
 =\sum_{i=0}^{\lfloor f/2\rfloor}
 C_{\lambda,f,i}(d)s_{f-i,i}(a,b),\qquad f\geq c.
\]
Here $Y_\lambda(d)$ is the \emph{open} coincident-root stratum.
With $\widetilde C_{\lambda,f,i}=(-1)^{f-c}C_{\lambda,f,i}$, the
assertion requires all coefficients $\gamma_j$ in
\[
 \widetilde C_{\lambda,f,i}\bigl(y+|\lambda|+2(f-c)\bigr)
       =\sum_{j\geq0}\gamma_j\binom yj
\]
to be nonnegative, with the resulting sequence log-concave.
As formulated, this includes the empty partition: it has no forbidden
parts, and the source explicitly discusses it immediately before the
conjecture. This case already contradicts the positivity requirement.

\begin{proposition}
For the empty partition, $f=1$, and $i=0$, the stable coefficient is
\[
 C_{\varnothing,1,0}(d)=\frac{d(3-d)}2.
\]
Consequently the conjecture's shifted signed polynomial is
\[
 \widetilde C_{\varnothing,1,0}(y+2)
   =\frac{(y+2)(y-1)}2
   =-\binom y0+\binom y1+\binom y2 .
\]
In particular, $\gamma_0=-1<0$.
\end{proposition}

\begin{proof}
Let $V_d=\operatorname{Pol}^d(\mathbb C^2)$, using the source's
positive-weight convention. Its weights are $ja+(d-j)b$ for
$0\leq j\leq d$, and therefore
\[
 c_1(V_d)=\sum_{j=0}^d(ja+(d-j)b)
          =\frac{d(d+1)}2(a+b).
\]
For every $d\geq2$, the open stratum $Y_\varnothing(d)$ is the
complement of the binary discriminant hypersurface $D_d$.
The discriminant has character $(\det)^{d(d-1)}$, so
\[
 [D_d]=d(d-1)(a+b).
\]
For completeness, write a split form as $\prod_{r=1}^d\ell_r$.
Its discriminant is, up to a nonzero constant,
$\prod_{r<s}\det(\ell_r,\ell_s)^2$. A linear change of variables
multiplies each determinant by $\det$, which proves the stated
character. The discriminant hypersurface is irreducible, being the
closure of the forms $\ell^2g$ with one double root. It is reduced:
along $(X^2-tY^2)g(X,Y)$ with the roots of $g$ distinct and away
from $X=0$, its discriminant is $t$ times a function nonzero at $t=0$.

The normalization and additivity of equivariant CSM classes give
\[
 \cSM(V_d)=c(V_d),\qquad
 \cSM(Y_\varnothing(d))=\cSM(V_d)-\cSM(D_d).
\]
The codimension-one part of $\cSM(D_d)$ is its fundamental class.
Thus
\[
 \cSM(Y_\varnothing(d))_1
   =c_1(V_d)-[D_d]=\frac{d(3-d)}2(a+b).
\]
Since $s_{1,0}(a,b)=a+b$, this identifies the Schur coefficient for
every $d\geq2$, and hence identifies its stable polynomial.
Now $c=0$, the sign is $-1$, and the prescribed shift is $y=d-2$.
Substitution gives the displayed binomial expansion.
Alternatively, evaluation at $y=0$ already forces its constant
binomial coefficient to be $-1$.
\end{proof}

This is a counterexample to the universal formulation, not a resolution
of the conjecture after imposing $\lambda\ne\varnothing$.
The distinction between the open stratum and its closure is essential:
$\overline{Y_\varnothing(d)}=V_d$, whereas $Y_\varnothing(d)$ consists
of squarefree forms. The source also distinguishes these objects after
Conjecture 6.7. No assertion about classes of closures is used here.

\paragraph{Task 2.}
Determine whether both conclusions of Conjecture 6.7 hold for every
nonempty partition $\lambda$ with parts at least $2$.

\paragraph{Task 3.}
For $\lambda=\varnothing$, determine the stable Schur coefficient
polynomials in all homogeneous degrees and the least integer shifts
that make their signed binomial coefficients nonnegative.

\paragraph{Task 4.}
For those corrected empty-partition shifts, determine which coefficient
sequences are log-concave and classify any exceptions.

\paragraph{Task 5.}
Formulate and test a separate binomial positivity and log-concavity
statement for CSM classes of coincident-root stratum closures, keeping
their normalization distinct from that of open strata.

\begin{thebibliography}{9}
\raggedright
\bibitem{BF}
G. B\'erczi and L. M. Feh\'er,
\emph{Positivity in classical enumerative geometry: a case study in
synchronized AI-assisted mathematics}, arXiv:2605.25271v1 (2026),
Section 6.2, pp.~26--27.
\url{https://arxiv.org/abs/2605.25271}.
\bibitem{J}
A. Juh\'asz, \emph{Equivariant classes of coincident root strata and
related invariants of varieties of tangent lines},
doctoral dissertation, Central European University (2023),
Chapter 3, in particular Theorem 3.6.1.
\url{https://hdl.handle.net/20.500.14018/29523}.
\end{thebibliography}
\end{document}
